% ==========================================================
%  OLD EXAM QUESTIONS — Reinforcement Learning
% ==========================================================

\documentclass[11pt,oneside]{article}

\usepackage[
  a4paper,
  top=2.4cm,
  bottom=2.1cm,
  left=2.3cm,
  right=2.3cm,
  headheight=15pt
]{geometry}

\usepackage[T1]{fontenc}
\usepackage[utf8]{inputenc}
\usepackage{lmodern}
\usepackage{microtype}
\usepackage{inconsolata}
\usepackage{amsmath,amssymb,mathtools,bm}
\usepackage{enumitem}
\usepackage{booktabs,array,multicol,longtable}
\usepackage{tabularx}
\usepackage{xcolor}
\usepackage[most]{tcolorbox}
\usepackage{titlesec}
\usepackage{fancyhdr}
\usepackage[hidelinks]{hyperref}

\setlength{\parindent}{0pt}
\setlength{\parskip}{0.55em}
\setlist[itemize]{topsep=3pt,itemsep=2pt,leftmargin=2.2em}
\setlist[enumerate]{topsep=3pt,itemsep=2pt,leftmargin=2.2em}

% ── Colors ───────────────────────────────────────────────
\definecolor{accent}{HTML}{1B3A5C}
\definecolor{q@frame}{HTML}{1A5276}
\definecolor{q@bg}{HTML}{F6FAFD}
\definecolor{key@frame}{HTML}{1E8449}
\definecolor{key@bg}{HTML}{F4FBF6}
\definecolor{note@frame}{HTML}{9A6A00}
\definecolor{note@bg}{HTML}{FEF9E7}

% ── Box global style ─────────────────────────────────────
\tcbset{
  enhanced,
  breakable,
  arc=2.5mm,
  boxrule=0.5pt,
  left=4.5mm,
  right=4mm,
  top=2.5mm,
  bottom=2.5mm,
  fonttitle=\bfseries\small,
  sharp corners=west,
  parbox=false
}

% Question box (blue) — one per exam question.
\newtcolorbox{qbox}[1][]{
  colback=q@bg,
  colframe=q@frame,
  borderline west={3.5pt}{0pt}{q@frame},
  title={#1}
}

% Key insight box (green) — answer hints, key results.
\newtcolorbox{keybox}[1][]{
  colback=key@bg,
  colframe=key@frame,
  borderline west={3.5pt}{0pt}{key@frame},
  title={#1}
}

% Organizational note box (amber) — structural remarks only.
\newtcolorbox{notebox}[1][]{
  colback=note@bg,
  colframe=note@frame,
  borderline west={3.5pt}{0pt}{note@frame},
  title={#1}
}

% ── Section headings ─────────────────────────────────────
\titleformat{\section}
  {\normalfont\Large\bfseries\color{accent}}
  {}{0em}{}
  [{\vspace{1pt}\color{accent!45}\titlerule[0.45pt]}]
\titlespacing*{\section}{0pt}{3.2ex plus 1ex minus .2ex}{1.8ex plus .2ex}

\titleformat{\subsection}
  {\normalfont\normalsize\bfseries\color{accent!80!black}}
  {}{0em}{}
\titlespacing*{\subsection}{0pt}{2.2ex plus 1ex minus .2ex}{1.0ex plus .2ex}

% ── Header & footer ──────────────────────────────────────
\pagestyle{fancy}
\fancyhf{}
\fancyhead[L]{\small\color{gray!70}Old Exam Questions — \coursetitle}
\fancyhead[R]{\small\color{gray!70}\thepage}
\renewcommand{\headrulewidth}{0.4pt}

% ── Document metadata ─────────────────────────────────────
\newcommand{\coursetitle}{Reinforcement Learning}

% ── Question counter & macros ────────────────────────────
\newcounter{question}
\newcommand{\questiontitle}[1]{%
  \stepcounter{question}\begin{qbox}[Question \thequestion: #1]}
\newcommand{\closequestion}{\end{qbox}}

% Mark questions that appeared on multiple exams.
\newcommand{\repeatnote}{\textit{\color{gray!70}Asked multiple times.}}

% Multiple-choice checkbox bullet.
\newcommand{\choice}{\item[$\square$] }

% Fill-in answer line. Usage: \answerline{3cm}
\newcommand{\answerline}[1]{\makebox[#1]{\hrulefill}}

% Blank working area — dashed-border box for student calculations.
\newcommand{\workingspace}[1]{%
  \par\vspace{2pt}%
  \begin{tcolorbox}[
    enhanced,
    breakable=false,
    colback=white,
    colframe=gray!40,
    boxrule=0.45pt,
    arc=0pt,
    sharp corners,
    height=#1,
    left=3mm, right=3mm, top=3mm, bottom=3mm,
    frame style={dashed}
  ]%
  \end{tcolorbox}%
  \vspace{2pt}%
}

% ── Math macros ──────────────────────────────────────────
\newcommand{\R}{\mathbb{R}}
\newcommand{\E}{\mathbb{E}}
\newcommand{\Prob}{\mathbb{P}}
\newcommand{\KL}[2]{D_{\mathrm{KL}}\!\left(#1\,\|\,#2\right)}
\newcommand{\dd}{\,\mathrm{d}}


% ==========================================================
%  BEGIN DOCUMENT
% ==========================================================
\begin{document}

% ── Title page ───────────────────────────────────────────
\thispagestyle{empty}
\vspace*{1.4cm}
\begin{center}
  {\color{accent}\rule{\linewidth}{1.8pt}}\par
  \vspace{0.65cm}
  {\fontsize{26}{32}\selectfont\bfseries\color{accent}Old Exam Questions\\[2pt] \coursetitle\par}
  \vspace{0.55cm}
  {\color{accent!35}\rule{\linewidth}{0.5pt}}\par
  \vspace{0.55cm}
  {\large Stefan Eder\par}
  \vspace{0.15cm}
  {\small\color{gray!70}\today\par}
\end{center}

\vspace{0.9cm}
\begin{notebox}[How this file is organized]
This file is separate from the course summary and is meant as a clean practice set.
Exact duplicates were removed. When essentially the same question appeared in multiple old exams,
that is marked with \repeatnote\ inside the question box.
The wording was lightly normalized where the scans were messy, but the mathematical content
and answer choices were kept intact.
\end{notebox}

\tableofcontents
\newpage


% ==========================================================
%  SECTION 1: Introduction & Basics
% ==========================================================

\section{Introduction \& Basics}

\questiontitle{Match Learning Paradigms}
Match the following definitions to the concept.
\begin{itemize}
    \item \textbf{Definition 1:} Learn a function that maps inputs (features) to outputs (targets), based on a training set consisting of example input-output pairs.
    \item \textbf{Definition 2:} Learn how agents should take actions in an environment in order to maximize a reward.
    \item \textbf{Definition 3:} Learn patterns and find structure from unlabeled data.
\end{itemize}
\textbf{Concepts:} Reinforcement Learning, Supervised Learning, Unsupervised Learning.

\medskip
Definition 1 $\rightarrow$ \answerline{4cm}\quad
Definition 2 $\rightarrow$ \answerline{4cm}\quad
Definition 3 $\rightarrow$ \answerline{4cm}
\closequestion

\questiontitle{Unsupervised Learning Definition}
Unsupervised learning considers the problem of goal-directed agents interacting with an uncertain environment.
\begin{itemize}[leftmargin=1.8em]
\choice True
\choice False
\end{itemize}
\closequestion

\questiontitle{RL Use Case}
Which of these applications is a good use case for reinforcement learning?
\begin{itemize}[leftmargin=1.8em]
\choice A software to identify which pictures show a cat, a dog, or a hamster in a photo library
\choice A computer program that plays chess
\choice Clustering a dataset of audio recordings to find similar songs
\choice An algorithm that identifies the different plants in your garden by taking a picture with your phone
\end{itemize}
\closequestion

\questiontitle{RL Components Match}
Match each definition with the concept. Each concept should be used only once.

\begin{tabularx}{\linewidth}{@{}>{\bfseries}l X >{\raggedleft\arraybackslash}p{4.2cm}@{}}
A. & Defines the agent's way of behaving at a given time & \answerline{4cm}\\[0.8ex]
B. & A model that mimics the behavior of the environment & \answerline{4cm}\\[0.8ex]
C. & A numerical value that defines what are good and bad events for the agent. The environment provides this value at each time step. & \answerline{4cm}\\[0.8ex]
D. & Indicates long-term desirability of states and is defined as the sum of all rewards that an agent can expect to accumulate over the future & \answerline{4cm}\\
\end{tabularx}
\textbf{Concepts:} reward signal, policy, model of the environment, value function.
\closequestion

\questiontitle{Environment, Agent, Policy Match}
Match each definition with the concept.

\begin{tabularx}{\linewidth}{@{}>{\bfseries}l X >{\raggedleft\arraybackslash}p{4.2cm}@{}}
A. & Setting or ``world'' within which the agent operates and interacts. & \answerline{4cm}\\[0.8ex]
B. & System or entity that makes decisions and takes actions based on the given policy, interacting with the environment. & \answerline{4cm}\\[0.8ex]
C. & Mapping from the state of the environment to the actions to be taken by the agent in those states. & \answerline{4cm}\\
\end{tabularx}
\textbf{Concepts:} environment; policy; agent.
\closequestion

\questiontitle{RL vs Supervised Learning}
Reinforcement learning is similar to supervised learning in which they both require explicit supervision in the form of a training set of inputs (features) and output (targets).
\begin{itemize}[leftmargin=1.8em]
\choice True
\choice False
\end{itemize}
\closequestion

\questiontitle{Exploration Definition}
In the context of the exploration-exploitation trade-off, what is one possible interpretation of the term exploration?
\begin{itemize}[leftmargin=1.8em]
\choice Always selecting actions randomly from a given probability distribution
\choice Taking actions that led to high rewards in the past
\choice Not always taking the action with the highest reward
\choice An agent must balance taking actions that have proven to be beneficial in the past and trying actions that have not been selected before
\end{itemize}
\closequestion

\questiontitle{SL/UL Exploration Trade-off}
Both unsupervised and supervised learning need to deal with the trade-off between exploration and exploitation.
\begin{itemize}[leftmargin=1.8em]
\choice True
\choice False
\end{itemize}
\closequestion


% ==========================================================
%  SECTION 2: Multi-Armed Bandits
% ==========================================================

\section{Multi-Armed Bandits}

\questiontitle{Exploration-Exploitation Methods in MAB}
Which of these methods explicitly (i.e., with a hyperparameter) balance exploration and exploitation in multi-armed bandits?
\begin{itemize}[leftmargin=1.8em]
\choice Upper Confidence Bound (UCB) action selection.
\choice Epsilon-greedy action selection
\choice Greedy action selection.
\choice The Gradient Bandit Algorithm
\end{itemize}
\closequestion

\questiontitle{Epsilon-Greedy Probability}
In epsilon greedy action selection, for the case of 10 actions and $\epsilon=0.7$, what is the probability that the greedy action is selected?
\begin{itemize}[leftmargin=1.8em]
\choice 0.37
\choice 0.7
\choice 0.03
\choice 0.1
\end{itemize}
\closequestion

\questiontitle{ISAM Value Update}
The arm $a$ of a slot machine has a true mean reward of $q(a)=-1$. After $t=10$ plays, your estimate of the mean reward is $Q_{t}(a)=-3$. The next play yields a reward of 2. What is the updated mean reward $Q_{t+1}$ using the Incremental Sample-Averaging Method (ISAM) and a learning rate $\alpha=0.1$?
\begin{itemize}[leftmargin=1.8em]
\choice $-2.5$
\choice $-2.9$
\choice $-3.1$
\choice $-0.5$
\end{itemize}
\closequestion

\questiontitle{k-Armed Bandit Sample Averages}
Consider a k-armed bandit problem with $k=2$ actions (1 and 2). We use $\epsilon$-greedy action selection and initial estimates $Q_1(1)=0$ and $Q_1(2)=0$. Suppose the initial sequence is:
\begin{itemize}
    \item $t=1$: $A_1=1, R_1=2$
    \item $t=2$: $A_2=1, R_2=-4$
    \item $t=3$: $A_3=2, R_3=-2$
    \item $t=4$: $A_4=1, R_4=-4$
    \item $t=5$: $A_5=2, R_5=-6$
    \item $t=6$: $A_6=2, R_6=-4$
\end{itemize}
Calculate the final value estimates $Q_7(1)$ and $Q_7(2)$ assuming sample-average updates.
\workingspace{4cm}
$Q_7(1) = $ \answerline{3cm}\qquad $Q_7(2) = $ \answerline{3cm}
\closequestion

\questiontitle{ISAM Step Size Role}
The step size $\alpha$ in the action-value update equation of ISAM \ldots
\begin{itemize}[leftmargin=1.8em]
\choice \ldots sets the probability of selecting a specific arm.
\choice \ldots adjusts the regret over time.
\choice \ldots controls how fast the algorithm explores new arms.
\choice \ldots determines the weight of new rewards in updating value estimates.
\end{itemize}
\closequestion

\questiontitle{Optimistic Initial Values Prior Knowledge}
Optimistic initial values requires some prior knowledge about the reward distribution.
\begin{itemize}[leftmargin=1.8em]
\choice True
\choice False
\end{itemize}
\closequestion

\questiontitle{Optimistic Initial Values Non-Stationary}
The optimistic initial values method is more effective in non-stationary environments than stationary ones.
\begin{itemize}[leftmargin=1.8em]
\choice True
\choice False
\end{itemize}
\closequestion

\questiontitle{UCB Uncertainty Decrease}
In UCB, the uncertainty term for an action that is not selected can decrease.
\begin{itemize}[leftmargin=1.8em]
\choice True
\choice False
\end{itemize}
\closequestion

\questiontitle{UCB Always Chooses Highest Value}
In the UCB algorithm, actions with higher estimated value $Q(a)$ are always chosen.
\begin{itemize}[leftmargin=1.8em]
\choice True
\choice False
\end{itemize}
\closequestion

\questiontitle{UCB Non-Stationary Suitability}
Upper Confidence Bound (UCB) algorithms are well-suited for non-stationary problems.
\begin{itemize}[leftmargin=1.8em]
\choice True
\choice False
\end{itemize}
\closequestion

\questiontitle{Exponential Recency-Weighted Averages}
Exponential recency-weighted averages can be used to handle non-stationary problems.
\begin{itemize}[leftmargin=1.8em]
\choice True
\choice False
\end{itemize}
\closequestion

\questiontitle{Gradient Bandit Model-Based}
The Gradient Bandit algorithm is a form of model-based reinforcement learning and requires an approximation of the reward distribution.
\begin{itemize}[leftmargin=1.8em]
\choice True
\choice False
\end{itemize}
\closequestion

\questiontitle{Gradient Bandit Action Selection}
In the Gradient Bandit algorithm, actions are selected proportional to their action selection preference.
\begin{itemize}[leftmargin=1.8em]
\choice True
\choice False
\end{itemize}
\closequestion

\questiontitle{Optimistic Initial Values Exploration Bias}
Using optimistic initial values biases the agent toward initially exploring actions with higher estimated rewards.
\begin{itemize}[leftmargin=1.8em]
\choice True
\choice False
\end{itemize}
\closequestion

\questiontitle{UCB Uncertainty Term}
The UCB algorithm balances exploration and exploitation by adding an uncertainty term to the estimated value of each action.
\begin{itemize}[leftmargin=1.8em]
\choice True
\choice False
\end{itemize}
\closequestion

\questiontitle{Non-Stationary Exploration Required}
Non-stationary problems do not require balancing exploration and exploitation.
\begin{itemize}[leftmargin=1.8em]
\choice True
\choice False
\end{itemize}
\closequestion

\questiontitle{Gradient Bandit Preference Updates}
In the Gradient Bandit algorithm, action selection preferences are updated after each action.
\begin{itemize}[leftmargin=1.8em]
\choice True
\choice False
\end{itemize}
\closequestion

\questiontitle{Epsilon Case First Occurrence}
On which time steps did the $\epsilon$-case definitely occur for the first time?
\begin{itemize}[leftmargin=1.8em]
\choice 2
\choice 5
\choice 6
\end{itemize}
\closequestion

\questiontitle{Best Asymptotic Method}
Which method will, on average, perform best in the long run (asymptotically) in terms of cumulative reward?
\begin{itemize}[leftmargin=1.8em]
\choice ISAM, $\epsilon = 0.1$
\choice ISAM, $\epsilon = 0.01$
\choice ISAM, greedy
\end{itemize}
\closequestion

\questiontitle{Optimistic Initial Values Negative Rewards}
If the rewards are strictly negative ($R < 0$), the Optimistic Initial Values method benefits from low initial value estimates ($Q_1 \ll 0$).
\begin{itemize}[leftmargin=1.8em]
\choice True
\choice False
\end{itemize}
\closequestion

\questiontitle{Optimistic Initial Values Non-Stationary Tracking}
Optimistic Initial Value methods are effective in tracking non-stationary k-armed bandit problems.
\begin{itemize}[leftmargin=1.8em]
\choice True
\choice False
\end{itemize}
\closequestion

\questiontitle{UCB with c=0}
UCB with $c = 0$ corresponds exactly to the greedy action selection method.
\begin{itemize}[leftmargin=1.8em]
\choice True
\choice False
\end{itemize}
\closequestion

\questiontitle{UCB Uses Epsilon-Greedy}
UCB relies on $\epsilon$-greedy action selection for exploration.
\begin{itemize}[leftmargin=1.8em]
\choice True
\choice False
\end{itemize}
\closequestion

\questiontitle{UCB Tries Each Action Once}
UCB initially selects each as-yet-untried action exactly once.
\begin{itemize}[leftmargin=1.8em]
\choice True
\choice False
\end{itemize}
\closequestion

\questiontitle{Gradient Bandit Random Exploration}
Gradient Bandit algorithms do not rely on random action selection for exploration.
\begin{itemize}[leftmargin=1.8em]
\choice True
\choice False
\end{itemize}
\closequestion

\questiontitle{Action Preferences vs Expected Reward}
Numerical action selection preferences cannot be translated to estimates of the expected reward.
\begin{itemize}[leftmargin=1.8em]
\choice True
\choice False
\end{itemize}
\closequestion

\questiontitle{Gradient Bandit Q-Value Baseline}
The gradient bandit algorithm can be improved by using action value estimates $Q(a)$ as the baseline in the update rule instead of $\bar{R}_t$.
\begin{itemize}[leftmargin=1.8em]
\choice True
\choice False
\end{itemize}
\closequestion


% ==========================================================
%  SECTION 3: Finite Markov Decision Processes
% ==========================================================

\section{Finite Markov Decision Processes}

\questiontitle{Stochastic Process Concepts Match}
Match each definition to the concept. Each concept should only be used once.

\begin{tabularx}{\linewidth}{@{}>{\bfseries}l X >{\raggedleft\arraybackslash}p{4.2cm}@{}}
A. & A collection of random variables $\{S_t \mid t \in T\}$. The variables take values from a state space $S$, and are indexed by some ordered index set $T$. & \answerline{4cm}\\[0.8ex]
B. & Property according to which a stochastic process can be classified as discrete-space or continuous-space. & \answerline{4cm}\\[0.8ex]
C. & The future is independent of the past given the present, i.e., $p(S_{t+1}\mid S_t, S_{t-1}, \ldots) = p(S_{t+1}\mid S_t)$. & \answerline{4cm}\\
\end{tabularx}
\textbf{Concepts:} Markov property; Stochastic Process; Cardinality of state space.
\closequestion

\questiontitle{MDP Equations Match}
Match the following definitions/equations about Markov decision processes (MDPs) with the concept.
\begin{itemize}
    \item $G_t = \sum_{k=0}^{\infty}\gamma^k R_{t+k+1}$ \hfill \answerline{4cm}
    \item $G_t = R_{t+1} + R_{t+2} + \cdots + R_T$ \hfill \answerline{4cm}
    \item $p(s', r|s, a) = \text{Pr}\{S_t = s', R_t = r | S_{t-1} = s, A_{t-1} = a\}$ \hfill \answerline{4cm}
\end{itemize}
\textbf{Concepts:} Discounted return; Return (undiscounted); Dynamics of the MDP.
\closequestion

\questiontitle{Discounted Return Calculation}
Suppose $\gamma = 1/2$ and the following sequence of rewards is received: $R_1 = -7, R_2 = 6, R_3 = 1, R_4 = 5, R_5 = 2$ with $T=5$ (assume $R_t=0$ for $t > T$). What are the returns $G_0, G_1, G_2, G_3, G_4, G_5$? (Hint: Work backwards.)
\workingspace{4cm}
$G_5 = $ \answerline{2cm},\quad
$G_4 = $ \answerline{2cm},\quad
$G_3 = $ \answerline{2cm},\quad
$G_2 = $ \answerline{2cm},\quad
$G_1 = $ \answerline{2cm},\quad
$G_0 = $ \answerline{2cm}
\closequestion

\questiontitle{Undiscounted Return Non-Episodic}
It is better to use the undiscounted return (i.e., $\gamma=1$) for non-episodic tasks.
\begin{itemize}[leftmargin=1.8em]
\choice True
\choice False
\end{itemize}
\closequestion

\questiontitle{MDP Action Influences Next State}
In general, in a Markov decision process, the decision of an action $a$ in state $S_t$ does not influence what the next state $S_{t+1}$ will be.
\begin{itemize}[leftmargin=1.8em]
\choice True
\choice False
\end{itemize}
\closequestion

\questiontitle{Infinite Discounted Return}
Consider a non-episodic task where there is a constant reward of $R=2$. What would be the discounted return after an infinite number of time steps, if the discount factor is $\gamma=0.75$?

$G_t = $ \answerline{3cm}
\closequestion


% ==========================================================
%  SECTION 4: Dynamic Programming
% ==========================================================

\section{Dynamic Programming}

\questiontitle{Policy Evaluation Computes Optimal Policy}
Policy evaluation computes the optimal policy directly.
\begin{itemize}[leftmargin=1.8em]
\choice True
\choice False
\end{itemize}
\closequestion

\questiontitle{Two-Array Policy Evaluation Update Order}
When using the two-array version of the policy evaluation algorithm, the order in which states are updated can affect the speed of convergence.
\begin{itemize}[leftmargin=1.8em]
\choice True
\choice False
\end{itemize}
\closequestion

\questiontitle{Greedy Policy in Policy Improvement}
In the context of policy improvement, `greedy policy' refers to a policy that explores actions randomly to gather more information.
\begin{itemize}[leftmargin=1.8em]
\choice True
\choice False
\end{itemize}
\closequestion

\questiontitle{One Policy Always Better}
Given two policies, one of them is guaranteed to be better than the other.
\begin{itemize}[leftmargin=1.8em]
\choice True
\choice False
\end{itemize}
\closequestion

\questiontitle{Optimal Policy Greedy w.r.t.\ Value}
The optimal policy is found if the policy is greedy with respect to its own value function.
\begin{itemize}[leftmargin=1.8em]
\choice True
\choice False
\end{itemize}
\closequestion

\questiontitle{Policy Iteration Finite State Space Only}
The policy iteration algorithm introduced in the lecture can be used for both finite and infinite state spaces.
\begin{itemize}[leftmargin=1.8em]
\choice True
\choice False
\end{itemize}
\closequestion

\questiontitle{Value Iteration Greedy Policy}
Value iteration indirectly evaluates a greedy policy with respect to the current value function estimate.
\begin{itemize}[leftmargin=1.8em]
\choice True
\choice False
\end{itemize}
\closequestion

\questiontitle{Value Iteration Alternates Eval and Improvement}
Value iteration alternates between policy evaluation and policy improvement.
\begin{itemize}[leftmargin=1.8em]
\choice True
\choice False
\end{itemize}
\closequestion

\questiontitle{Policy Evaluation Convergence Guarantee}
Policy evaluation is guaranteed to converge to the value function of the given policy.
\begin{itemize}[leftmargin=1.8em]
\choice True
\choice False
\end{itemize}
\closequestion

\questiontitle{Policy Improvement Strictly Better}
Policy improvement in dynamic programming always leads to a strictly better policy.
\begin{itemize}[leftmargin=1.8em]
\choice True
\choice False
\end{itemize}
\closequestion

\questiontitle{Policy Iteration Not Guaranteed Finite}
In policy iteration, an optimal policy is not guaranteed to be found after a finite number of steps.
\begin{itemize}[leftmargin=1.8em]
\choice True
\choice False
\end{itemize}
\closequestion

\questiontitle{Policy Iteration Guaranteed Finite}
\repeatnote
Policy Iteration is guaranteed to find an optimal policy after a finite number of steps.
\begin{itemize}[leftmargin=1.8em]
\choice True
\choice False
\end{itemize}
\closequestion

\questiontitle{Iterative Policy Evaluation Optimal Policy}
Iterative Policy Evaluation is guaranteed to find the optimal policy after a finite number of steps.
\begin{itemize}[leftmargin=1.8em]
\choice True
\choice False
\end{itemize}
\closequestion

\questiontitle{Value Iteration Finite Convergence}
Value iteration converges to the optimal value function after a finite number of iterations.
\begin{itemize}[leftmargin=1.8em]
\choice True
\choice False
\end{itemize}
\closequestion

\questiontitle{Value Iteration Stochastic Unique Optimum}
If the environment is stochastic, value iteration always converges to a unique optimal value function.
\begin{itemize}[leftmargin=1.8em]
\choice True
\choice False
\end{itemize}
\closequestion

\questiontitle{Value Iteration Bellman Optimality Update}
Value iteration uses a Bellman-optimality-equation-like update rule.
\begin{itemize}[leftmargin=1.8em]
\choice True
\choice False
\end{itemize}
\closequestion

\questiontitle{Iterative Policy Evaluation Bootstrapping}
Iterative Policy Evaluation uses bootstrapping to estimate a state's value.
\begin{itemize}[leftmargin=1.8em]
\choice True
\choice False
\end{itemize}
\closequestion

\questiontitle{Iterative Policy Evaluation Successor States}
Iterative Policy Evaluation uses the value estimates of all immediate successor states to update the current state's value.
\begin{itemize}[leftmargin=1.8em]
\choice True
\choice False
\end{itemize}
\closequestion

\questiontitle{Optimal Policy via Linear Equations}
Given an arbitrary policy, the optimal policy can be computed directly by solving a system of linear equations using the Bellman equation.
\begin{itemize}[leftmargin=1.8em]
\choice True
\choice False
\end{itemize}
\closequestion

\questiontitle{Policy Improvement Small vs Large Step}
Consider a state 0 where the agent can choose between taking a \textbf{small} or a \textbf{large} step to the right.
\begin{itemize}
    \item The \textbf{small step} takes the agent into state 1 or state 3 with probability 0.5 each.
    \item The \textbf{large step} takes the agent into state 2 (prob 0.5), state 4 (prob 0.25), or state 5 (prob 0.25).
\end{itemize}
Given the following values: $V(1)=3, V(2)=5, V(3)=-1, V(4)=-3, V(5)=-7$, reward $R=-1$ for all transitions, and $\gamma=1$.
\workingspace{4cm}
\begin{enumerate}
    \item Calculate $Q(\text{small})$: \answerline{3cm}
    \item Calculate $Q(\text{large})$: \answerline{3cm}
    \item Which action does policy improvement select? \answerline{3cm}
\end{enumerate}
\closequestion

\questiontitle{DP Iterative Policy Evaluation Grid}
Consider a grid world with states 1, 2, and 3, plus a terminal state (grey). Reward $= -1$ per transition. Deterministic environment. Discount $\gamma = 0.5$.

Current values: $V(1)=0, V(2)=0, V(3)=0$.

Perform one sweep of iterative policy evaluation (Two-array version).
\workingspace{4cm}
\begin{itemize}
    \item $V_{t+1}(1)$: \answerline{3cm}
    \item $V_{t+1}(2)$: \answerline{3cm}
    \item $V_{t+1}(3)$: \answerline{3cm}
\end{itemize}
\closequestion

\questiontitle{Policy Improvement Grid Q-Values}
Consider a grid world with the following current value estimates $V_t(s)$:
\begin{center}
\begin{tabular}{|c|c|c|}
\hline
State 1 ($V=4$) & State 2 ($V=2$) & State 3 ($V=0$) \\
\hline
 & State 4 ($V=8$) & \\
\hline
\end{tabular}
\end{center}
The agent receives a reward of $R=-1$ for each transition. The environment is stochastic: moves succeed with probability 0.5, or slide to a perpendicular cell with probability 0.25 each. $\gamma=1.0$.

Calculate the q-values for State 2:
\workingspace{5cm}
\begin{itemize}
    \item $Q(2, \leftarrow)$: \answerline{3cm}
    \item $Q(2, \rightarrow)$: \answerline{3cm}
    \item $Q(2, \uparrow)$ (wall, assume stays put or slides): \answerline{3cm}
    \item $Q(2, \downarrow)$: \answerline{3cm}
\end{itemize}
\closequestion

\questiontitle{DP Grid Policy Improvement Action}
Which action will policy improvement select for state 2 in the Dynamic Programming Iterative Policy Evaluation (Grid) problem?
\begin{itemize}[leftmargin=1.8em]
\choice $\rightarrow$
\choice $\downarrow$
\choice $\uparrow$
\choice $\leftarrow$
\end{itemize}
\closequestion

\questiontitle{Policy Evaluation Two-Sweep Value}
Consider a grid world where the environment gives a reward of $-4$ for every transition and the discount factor is 0.5. Assume the terminal state value is always 0. What is the estimated value of state 0 after \textbf{two sweeps} of policy evaluation? (Assume initial values are 0.)

\answerline{3cm}
\closequestion


% ==========================================================
%  SECTION 5: Monte Carlo Methods
% ==========================================================

\section{Monte Carlo Methods}

\questiontitle{MC Control Methods Match}
For each of the following statements, select the corresponding method from the list below:
\begin{enumerate}
    \item \textbf{Statement A:} The probability of selecting the greedy action for the $k$-th episode is given by $\epsilon = 1/k$.
    \item \textbf{Statement B:} The probability of selecting the greedy action is given by $1-\epsilon$, with $\epsilon$ being a constant.
    \item \textbf{Statement C:} Learn about target policy $\pi$ from experience sampled from a behavior policy $\mu$.
\end{enumerate}
\textbf{Methods:}
\begin{itemize}[leftmargin=1.8em]
\choice Off-policy learning
\choice On-policy MC control for epsilon soft policies
\choice Greedy in the limit with infinite exploration (GLIE) MC control
\end{itemize}

Statement A $\rightarrow$ \answerline{5cm}

Statement B $\rightarrow$ \answerline{5cm}

Statement C $\rightarrow$ \answerline{5cm}
\closequestion

\questiontitle{MC Exploring Starts vs GLIE vs Off-Policy}
For each of the following statements, select the corresponding method:
\begin{enumerate}
    \item An $\epsilon$-greedy method where the value of $\epsilon$ decreases every episode. \hfill \answerline{4cm}
    \item Use importance sampling to weight the return according to the similarity between the target ($\pi$) and behavior ($\mu$) policies. \hfill \answerline{4cm}
    \item In this method, start each episode at random state-action pairs (all state-action pairs have non-zero probabilities of being at the start). \hfill \answerline{4cm}
\end{enumerate}
\textbf{Concepts:} GLIE; MC control with exploring starts; Off-policy learning.
\closequestion

\questiontitle{First-Visit MC Correlation}
Value estimates in the first-visit Monte Carlo Policy estimation algorithm are correlated with each other.
\begin{itemize}[leftmargin=1.8em]
\choice True
\choice False
\end{itemize}
\closequestion

\questiontitle{MC Control Estimates State Values}
Monte Carlo control methods compute an estimate of the state value function $v_\pi$.
\begin{itemize}[leftmargin=1.8em]
\choice True
\choice False
\end{itemize}
\closequestion

\questiontitle{MC Convergence Speed}
In practice, Monte Carlo methods for value estimation require just a few episodes to converge to the value function.
\begin{itemize}[leftmargin=1.8em]
\choice True
\choice False
\end{itemize}
\closequestion

\questiontitle{MC Requires Dynamics Knowledge}
Monte Carlo methods require explicit knowledge of the dynamics of the problem ($p(s',r|s,a)$).
\begin{itemize}[leftmargin=1.8em]
\choice True
\choice False
\end{itemize}
\closequestion

\questiontitle{Fish Lake First-Visit MC}
Following policy $\pi$, we sample the following 4 episodes in a ``Fish Lake'' grid world (actions: left, right; rewards: 0, 1, $-1$):
\begin{itemize}
    \item \textbf{Ep 1:} $S_0=b, A_0=\text{left}, R_1=0, S_1=b, A_1=\text{left}, R_2=0, S_2=b, A_2=\text{right}, R_3=1, S_3=c$
    \item \textbf{Ep 2:} $S_0=b, A_0=\text{left}, R_1=0, S_1=b, A_1=\text{left}, R_2=0, S_2=b, A_2=\text{left}, R_3=-1, S_3=a$
    \item \textbf{Ep 3:} $S_0=b, A_0=\text{right}, R_1=0, S_1=b, A_1=\text{right}, R_2=0, S_2=b, A_2=\text{right}, R_3=1, S_3=c$
    \item \textbf{Ep 4:} $S_0=b, A_0=\text{left}, R_1=0, S_1=b, A_1=\text{right}, R_2=0, S_2=b, A_2=\text{right}, R_3=1, S_3=c$
\end{itemize}
Using \textbf{First Visit Monte Carlo Policy Evaluation} ($\gamma=1$):
\workingspace{5cm}
\begin{enumerate}
    \item Estimate $q_\pi(b, \text{left})$: \answerline{3cm}
    \item Estimate $q_\pi(b, \text{right})$: \answerline{3cm}
\end{enumerate}
\closequestion

\questiontitle{MC Model-Free}
Monte Carlo methods are model-free.
\begin{itemize}[leftmargin=1.8em]
\choice True
\choice False
\end{itemize}
\closequestion

\questiontitle{Every-Visit MC IID Samples}
Samples from the every-visit Monte Carlo Policy estimation algorithm are independent and identically distributed.
\begin{itemize}[leftmargin=1.8em]
\choice True
\choice False
\end{itemize}
\closequestion


% ==========================================================
%  SECTION 6: Temporal Difference Learning
% ==========================================================

\section{Temporal Difference Learning}

\questiontitle{TD Policy Evaluation Grid S0-Start}
Consider a grid world where the agent follows a deterministic policy $\pi$:
\begin{itemize}
    \item From $S_0$, go Right to $S_1$.
    \item From $S_1$, go Up to $S_2$.
    \item From $S_2$, go Left to $S_T$ (Terminal).
\end{itemize}
Parameters: $\gamma = 1$, Reward $R = -1$ per step, $\alpha = 1$. Initial values $V(S) = 0$.

Perform TD(0) updates for one episode: $S_0 \to S_1 \to S_2 \to S_T$.

Calculate the updated values (round to 3 decimal places):
\workingspace{4cm}
\begin{itemize}
    \item $V(S_1)$ after the step $S_0 \to S_1$: \answerline{3cm}
    \item $V(S_2)$ after the step $S_1 \to S_2$: \answerline{3cm}
\end{itemize}
\closequestion

\questiontitle{Off-Policy SARSA Statements}
Assuming the MDP above (TD Policy Evaluation S0-Start), we use an off-policy version of SARSA, using the deterministic target policy $\pi$ as indicated by the arrows, and a non-deterministic behavior policy $\mu$ to explore. $Q(S_t, A_t)$ is initialized as 0 for all states and actions. The agent always starts the episode in state $S_0$. Which of the following statements are true?
\begin{itemize}[leftmargin=1.8em]
\choice Following the deterministic policy $\pi(a|S)$, some state action pairs are never explored.
\choice $Q(S_t, A_t)$ for actions $A_t$ chosen by the deterministic policy converges to the TD(0) policy evaluation solution $V(S_t)$ computed before.
\choice This version of SARSA has no optimality guarantees because the discount factor $\gamma$ is larger than 0.
\choice This off-policy version of SARSA would be identical to Q-learning.
\end{itemize}
\closequestion

\questiontitle{Q-Learning Off-Policy No IS}
Q-Learning is an off-policy method that does not require importance sampling.
\begin{itemize}[leftmargin=1.8em]
\choice True
\choice False
\end{itemize}
\closequestion

\questiontitle{TD Advantage over MC}
What is an advantage of TD learning over Monte Carlo methods?
\begin{itemize}[leftmargin=1.8em]
\choice In practice, TD learning converges faster because it doesn't need to wait for an episode to end.
\choice It always gives a more accurate estimate.
\choice TD learning requires the dynamics of the environment, while MC methods are model-free.
\choice It works better in deterministic environments.
\end{itemize}
\closequestion

\questiontitle{Expected SARSA Equals Q-Learning}
Expected SARSA and Q-Learning can be the same for some choice of policy.
\begin{itemize}[leftmargin=1.8em]
\choice True
\choice False
\end{itemize}
\closequestion

\questiontitle{TD Non-Episodic Tasks}
What is the primary reason TD learning can handle non-episodic tasks better than Monte Carlo methods?
\begin{itemize}[leftmargin=1.8em]
\choice TD learning does not require a discount factor.
\choice TD learning is model-free, while Monte Carlo is model-based.
\choice TD learning updates based on partial episodes, allowing for continuous learning.
\choice TD learning can compute exact state values in one step.
\end{itemize}
\closequestion

\questiontitle{SARSA Second Action Same Policy}
SARSA requires choosing a second action in the second state with the same policy as the first.
\begin{itemize}[leftmargin=1.8em]
\choice True
\choice False
\end{itemize}
\closequestion

\questiontitle{SARSA vs Expected SARSA Direction}
SARSA moves deterministically in the direction Expected-SARSA is moving in expectation.
\begin{itemize}[leftmargin=1.8em]
\choice True
\choice False
\end{itemize}
\closequestion

\questiontitle{N-Step TD Equivalent to Exhaustive Search}
TD methods can take many steps at the same time, effectively interpolating between TD(0) and Exhaustive Search methods.
\begin{itemize}[leftmargin=1.8em]
\choice True
\choice False
\end{itemize}
\closequestion

\questiontitle{N-Step Intermediate Performance}
N-step methods somewhere in between TD(0) and TD(1) will often perform better than either extreme.
\begin{itemize}[leftmargin=1.8em]
\choice True
\choice False
\end{itemize}
\closequestion

\questiontitle{TD(lambda) Weighted Average}
The TD($\lambda$) algorithm computes a weighted average n-step TD backup target.
\begin{itemize}[leftmargin=1.8em]
\choice True
\choice False
\end{itemize}
\closequestion

\questiontitle{Eligibility Traces Forward View}
Eligibility Traces implement a forward view of TD($\lambda$) methods.
\begin{itemize}[leftmargin=1.8em]
\choice True
\choice False
\end{itemize}
\closequestion

\questiontitle{TD(0) Update Rule}
Which of the following equations represents the TD(0) update rule for policy evaluation?
\begin{itemize}[leftmargin=1.8em]
\choice $V(s) \leftarrow V(s) + \alpha [G_t - V(s)]$
\choice $V(s) \leftarrow V(s) + \alpha [R + \gamma V(s') - V(s)]$
\choice $Q(s, a) \leftarrow Q(s, a) + \alpha [R + \gamma \max_{a'} Q(s', a') - Q(s, a)]$
\end{itemize}
\closequestion

\questiontitle{TD(0) Discount Factor Zero}
In policy evaluation using TD(0), what happens if the discount factor ($\gamma$) is set to 0?
\begin{itemize}[leftmargin=1.8em]
\choice The agent considers only immediate rewards.
\choice The agent evaluates the policy using the average reward per time step.
\choice The learning process diverges.
\choice The value function estimates future rewards up to an infinite horizon.
\end{itemize}
\closequestion

\questiontitle{TD(0) Bias vs MC}
TD(0) generally is less biased than Monte Carlo methods.
\begin{itemize}[leftmargin=1.8em]
\choice True
\choice False
\end{itemize}
\closequestion

\questiontitle{TD(0) vs MC Key Difference}
What is the key difference between TD(0) and Monte Carlo methods in policy evaluation?
\begin{itemize}[leftmargin=1.8em]
\choice TD(0) only uses undiscounted returns, while Monte Carlo uses both undiscounted and discounted returns.
\choice TD(0) requires a model of the environment; Monte Carlo does not.
\choice TD(0) updates after observing a full episode, while Monte Carlo updates after every step.
\choice TD(0) updates after every step, while Monte Carlo updates after observing a full episode.
\end{itemize}
\closequestion

\questiontitle{N-Step TD Equivalent to Dynamic Programming}
TD methods can take many steps at the same time. If the number of steps is greater than the episode length, TD methods are equivalent to Dynamic Programming.
\begin{itemize}[leftmargin=1.8em]
\choice True
\choice False
\end{itemize}
\closequestion

\questiontitle{N-Step TD Incomplete Episodes}
N-step TD methods have to wait $n$ time steps before making an update but can update from incomplete episodes.
\begin{itemize}[leftmargin=1.8em]
\choice True
\choice False
\end{itemize}
\closequestion

\questiontitle{TD Policy Evaluation S3-Start}
Consider a grid world with states 1, 2, 3 and Terminal. Reward $= -1$. Deterministic. $\gamma=0.5, \alpha=0.5$.

Current values: $V(1)=0, V(2)=0, V(3)=0$.

Episode: Start at $S_3 \to \text{Right} \to S_2 \to \text{Up} \to S_1 \to \text{Left} \to S_T$.

Calculate the updates:
\workingspace{4cm}
\begin{itemize}
    \item $V(S_3)$ after first step: \answerline{3cm}
    \item $V(S_2)$ after second step: \answerline{3cm}
\end{itemize}
\closequestion

\questiontitle{TD(0) Convergence Settings}
Assuming the MDP from the ``TD Policy Evaluation (S3 start)'' question, under which settings for the step size $\alpha$ and the discount rate $\gamma$ does the TD(0) estimate converge to the same estimate as the Dynamic Programming solution?
\begin{itemize}[leftmargin=1.8em]
\choice $\alpha = 0.1, \gamma = 0.9$
\choice $\alpha = 1.0, \gamma = 0.5$
\choice $\alpha = 0.5, \gamma = 0.5$
\choice $\alpha = 0.5, \gamma = 0.9$
\end{itemize}
\closequestion

\questiontitle{SARSA on S3-Start Grid}
Assuming the MDP from the ``TD Policy Evaluation (S3 start)'' question, we use the SARSA algorithm for TD(0) control while keeping the deterministic behavior policy. $Q(S_t, A_t)$ is initialized as in the figure for all actions. Which statements are true?
\begin{itemize}[leftmargin=1.8em]
\choice This version of SARSA has no optimality guarantees because the discount factor $\gamma$ is less than 1.0.
\choice This version of SARSA converges to the globally optimal policy.
\choice This version of SARSA has no optimality guarantees because some state action pairs are never explored.
\choice $Q(S_t, A_t)$ for actions chosen by the deterministic policy converges to the TD(0) policy evaluation solution computed before.
\end{itemize}
\closequestion

\questiontitle{TD(0) Exact Same Sequence as DP}
Under which settings for the step size $\alpha$ and the discount rate $\gamma$ does the TD(0) policy evaluation produce the \textbf{exact same sequence} of estimates as the Dynamic Programming in the earlier section of the exam?
\begin{itemize}[leftmargin=1.8em]
\choice $\alpha = 0.5, \gamma = 0.9$
\choice $\alpha = 0.5, \gamma = 1.0$
\choice $\alpha = 0.5, \gamma = 0.5$
\choice $\alpha = 1.0, \gamma = 0.5$
\end{itemize}
\closequestion

\questiontitle{Expected SARSA vs SARSA Non-Greedy}
Expected SARSA and SARSA can be the same for some choice of non-greedy behavior policy.
\begin{itemize}[leftmargin=1.8em]
\choice True
\choice False
\end{itemize}
\closequestion

\questiontitle{TD Requires Environment Model}
TD methods require a model of the environment (i.e., the dynamics and reward function).
\begin{itemize}[leftmargin=1.8em]
\choice True
\choice False
\end{itemize}
\closequestion

\questiontitle{TD(0) Bootstrapping}
TD(0) uses bootstrapping (i.e., they update the value function using previous estimates of the value function).
\begin{itemize}[leftmargin=1.8em]
\choice True
\choice False
\end{itemize}
\closequestion

\questiontitle{Batch MC vs Batch TD Maximum Likelihood}
Both batch Monte Carlo learning and batch Temporal Difference learning converge to the maximum likelihood Markov model.
\begin{itemize}[leftmargin=1.8em]
\choice True
\choice False
\end{itemize}
\closequestion

\questiontitle{N-Step TD Equivalent to MC}
TD methods can take many steps at the same time. If the number of steps is greater than the episode length, TD methods are equivalent to Monte Carlo methods.
\begin{itemize}[leftmargin=1.8em]
\choice True
\choice False
\end{itemize}
\closequestion

\questiontitle{TD(lambda) Semi-Gradient Convergence}
TD($\lambda$) targets can be used in semi-gradient methods without convergence guarantees.
\begin{itemize}[leftmargin=1.8em]
\choice True
\choice False
\end{itemize}
\closequestion


% ==========================================================
%  SECTION 7: Policy Gradient Methods
% ==========================================================

\section{Policy Gradient Methods}

\questiontitle{Policy Gradient Log Probability}
In the policy gradient theorem, the term $\nabla_\theta \log \pi_\theta(a|s)$ represents the direction in parameter space that increases the probability of selecting action $a$.
\begin{itemize}[leftmargin=1.8em]
\choice True
\choice False
\end{itemize}
\closequestion

\questiontitle{Policy Gradient Deterministic Only}
The policy gradient theorem applies only to deterministic policies.
\begin{itemize}[leftmargin=1.8em]
\choice True
\choice False
\end{itemize}
\closequestion

\questiontitle{REINFORCE Deterministic Only}
REINFORCE can only be used to learn deterministic policies.
\begin{itemize}[leftmargin=1.8em]
\choice True
\choice False
\end{itemize}
\closequestion

\questiontitle{REINFORCE Episodic Only}
REINFORCE is effective for episodic and non-episodic Markov decision processes.
\begin{itemize}[leftmargin=1.8em]
\choice True
\choice False
\end{itemize}
\closequestion

\questiontitle{REINFORCE Baseline Reduces Variance}
Using an appropriate baseline in REINFORCE with baseline reduces the variance of the gradient estimator.
\begin{itemize}[leftmargin=1.8em]
\choice True
\choice False
\end{itemize}
\closequestion

\questiontitle{REINFORCE with Baseline Actor-Critic}
The REINFORCE algorithm with baseline is considered an actor-critic method.
\begin{itemize}[leftmargin=1.8em]
\choice True
\choice False
\end{itemize}
\closequestion

\questiontitle{Actor-Critic Variance Issues}
Actor-critic methods suffer from the same variance issues as REINFORCE because they do not use baselines.
\begin{itemize}[leftmargin=1.8em]
\choice True
\choice False
\end{itemize}
\closequestion

\questiontitle{Actor-Critic Critic Gradient}
Actor-critic methods require the gradient of the critic to compute the policy update.
\begin{itemize}[leftmargin=1.8em]
\choice True
\choice False
\end{itemize}
\closequestion

\questiontitle{Policy Gradient State Distribution}
The policy gradient theorem states that the gradient of the expected reward can be expressed without the need to compute the gradient of the state distribution.
\begin{itemize}[leftmargin=1.8em]
\choice True
\choice False
\end{itemize}
\closequestion

\questiontitle{REINFORCE Requires Value Function}
REINFORCE requires a value function approximation to compute the policy update.
\begin{itemize}[leftmargin=1.8em]
\choice True
\choice False
\end{itemize}
\closequestion

\questiontitle{REINFORCE On-Policy}
The REINFORCE algorithm is an on-policy method, i.e., it requires trajectories generated by the current policy.
\begin{itemize}[leftmargin=1.8em]
\choice True
\choice False
\end{itemize}
\closequestion

\questiontitle{REINFORCE Baseline Must Be Constant}
The baseline in REINFORCE with baseline must be a constant value for the gradient estimate to remain unbiased.
\begin{itemize}[leftmargin=1.8em]
\choice True
\choice False
\end{itemize}
\closequestion

\questiontitle{REINFORCE Baseline as Value Function}
A common choice for the baseline in REINFORCE with baseline is an estimate of the value function $V(s)$.
\begin{itemize}[leftmargin=1.8em]
\choice True
\choice False
\end{itemize}
\closequestion

\questiontitle{Actor-Critic Critic Learns State Value}
The critic in actor-critic methods typically learns an approximation of the state value function $V(s)$.
\begin{itemize}[leftmargin=1.8em]
\choice True
\choice False
\end{itemize}
\closequestion

\questiontitle{PPO Clipped Surrogate Objective}
Proximal Policy Optimization (PPO) is an actor-critic method that uses a clipped surrogate objective to stabilize training.
\begin{itemize}[leftmargin=1.8em]
\choice True
\choice False
\end{itemize}
\closequestion

\questiontitle{Policy-Based RL Uses Action Values}
Policy-based RL methods rely on action values to select actions.
\begin{itemize}[leftmargin=1.8em]
\choice True
\choice False
\end{itemize}
\closequestion

\questiontitle{Supervised Pretraining for Policy RL}
Supervised pretraining with human-generated state-action pairs is an appropriate strategy to inject expert knowledge into a policy-based RL model.
\begin{itemize}[leftmargin=1.8em]
\choice True
\choice False
\end{itemize}
\closequestion

\questiontitle{Policy-Based RL Exploration via Sampling}
Policy-based RL methods trade-off exploration and exploitation by sampling from the parameterized policy.
\begin{itemize}[leftmargin=1.8em]
\choice True
\choice False
\end{itemize}
\closequestion

\questiontitle{Performance Gradient State Distribution}
The exact performance gradient $\nabla J(\Theta)$ depends on changes in the state distribution; a proportional quantity is often used as a substitute.
\begin{itemize}[leftmargin=1.8em]
\choice True
\choice False
\end{itemize}
\closequestion

\questiontitle{REINFORCE with Baseline TD Method}
The REINFORCE algorithm with baseline is considered a Temporal Difference learning method because it uses bootstrapping with state value estimates to estimate returns.
\begin{itemize}[leftmargin=1.8em]
\choice True
\choice False
\end{itemize}
\closequestion

\questiontitle{REINFORCE with Baseline Non-Episodic}
The REINFORCE algorithm with baseline can be used for non-episodic tasks.
\begin{itemize}[leftmargin=1.8em]
\choice True
\choice False
\end{itemize}
\closequestion


% ==========================================================
%  ANSWER KEY
% ==========================================================
\newpage
\section*{Answer Key}

\begin{enumerate}[label=\textbf{Q\arabic*.}, leftmargin=3em, itemsep=4pt]

  % Section 1: Introduction & Basics
  \item Definition 1 $\rightarrow$ Supervised Learning;\quad Definition 2 $\rightarrow$ Reinforcement Learning;\quad Definition 3 $\rightarrow$ Unsupervised Learning
  \item \textbf{False}
  \item A computer program that plays chess
  \item A $\rightarrow$ policy;\quad B $\rightarrow$ model of the environment;\quad C $\rightarrow$ reward signal;\quad D $\rightarrow$ value function
  \item A $\rightarrow$ environment;\quad B $\rightarrow$ agent;\quad C $\rightarrow$ policy
  \item \textbf{False}
  \item Not always taking the action with the highest reward
  \item \textbf{False}

  % Section 2: Multi-Armed Bandits
  \item UCB and Epsilon-greedy (both correct)
  \item 0.37
  \item $-2.5$
  \item $Q_7(1) = -2$;\quad $Q_7(2) = -4$
  \item \ldots determines the weight of new rewards in updating value estimates.
  \item \textbf{True}
  \item \textbf{False}
  \item \textbf{False}
  \item \textbf{False}
  \item \textbf{False}
  \item \textbf{True}
  \item \textbf{False}
  \item \textbf{True}
  \item \textbf{True}
  \item \textbf{False}
  \item \textbf{True}
  \item 5
  \item ISAM, $\epsilon = 0.01$
  \item \textbf{False}
  \item \textbf{False}
  \item \textbf{True}
  \item \textbf{False}
  \item \textbf{True}
  \item \textbf{False}

  % Section 3: Finite MDPs
  \item A $\rightarrow$ Stochastic Process;\quad B $\rightarrow$ Cardinality of state space;\quad C $\rightarrow$ Markov property
  \item First bullet $\rightarrow$ Discounted return;\quad Second bullet $\rightarrow$ Return (undiscounted);\quad Third bullet $\rightarrow$ Dynamics of the MDP
  \item $G_5=0$,\quad $G_4=2$,\quad $G_3=6$,\quad $G_2=4$,\quad $G_1=8$,\quad $G_0=-3$
  \item \textbf{False}
  \item \textbf{False}
  \item $G_t = 8$

  % Section 4: Dynamic Programming
  \item \textbf{False}
  \item \textbf{False}
  \item \textbf{False}
  \item \textbf{False}
  \item \textbf{True}
  \item \textbf{False}
  \item \textbf{True}
  \item \textbf{False}
  \item \textbf{True}
  \item \textbf{False}
  \item \textbf{False}
  \item \textbf{True}
  \item \textbf{False}
  \item \textbf{False}
  \item \textbf{True}
  \item \textbf{True}
  \item \textbf{True}
  \item \textbf{False}
  \item $Q(\text{small}) = 0$;\quad $Q(\text{large}) = -1$;\quad selects \textbf{small step}
  \item Requires grid figure from slides — see problem setup (values depend on neighbor positions)
  \item $\downarrow$
  \item $-6$ (after two sweeps: first sweep gives $-4$, second sweep gives $-4 + 0.5 \times (-4) = -6$)

  % Section 5: Monte Carlo Methods
  \item Statement A $\rightarrow$ GLIE MC control;\quad Statement B $\rightarrow$ On-policy MC control for epsilon soft policies;\quad Statement C $\rightarrow$ Off-policy learning
  \item 1 $\rightarrow$ GLIE;\quad 2 $\rightarrow$ Off-policy learning;\quad 3 $\rightarrow$ MC control with exploring starts
  \item \textbf{False}
  \item \textbf{False}
  \item \textbf{False}
  \item \textbf{False}
  \item $q_\pi(b, \text{left}) \approx 1/3 \approx 0.333$;\quad $q_\pi(b, \text{right}) = 1$
  \item \textbf{True}
  \item \textbf{False}

  % Section 6: Temporal Difference Learning
  \item $V(S_1) = -1$;\quad $V(S_2) = -1$
  \item Statements 1 and 2 are \textbf{True}; statements 3 and 4 are \textbf{False}
  \item \textbf{True}
  \item In practice, TD learning converges faster because it doesn't need to wait for an episode to end.
  \item \textbf{True}
  \item TD learning updates based on partial episodes, allowing for continuous learning.
  \item \textbf{True}
  \item \textbf{False}
  \item \textbf{False}
  \item \textbf{True}
  \item \textbf{True}
  \item \textbf{False}
  \item $V(s) \leftarrow V(s) + \alpha [R + \gamma V(s') - V(s)]$
  \item The agent considers only immediate rewards.
  \item \textbf{False}
  \item TD(0) updates after every step, while Monte Carlo updates after observing a full episode.
  \item \textbf{False}
  \item \textbf{True}
  \item $V(S_3) = -0.5$;\quad $V(S_2) = -0.5$
  \item $\alpha = 1.0, \gamma = 0.5$ and $\alpha = 0.5, \gamma = 0.5$ (both correct)
  \item Statements 3 and 4 are \textbf{True}; statements 1 and 2 are \textbf{False}
  \item $\alpha = 1.0, \gamma = 0.5$
  \item \textbf{False}
  \item \textbf{False}
  \item \textbf{True}
  \item \textbf{False}
  \item \textbf{True}
  \item \textbf{True}

  % Section 7: Policy Gradient Methods
  \item \textbf{True}
  \item \textbf{False}
  \item \textbf{False}
  \item \textbf{False}
  \item \textbf{True}
  \item \textbf{False}
  \item \textbf{False}
  \item \textbf{False}
  \item \textbf{True}
  \item \textbf{False}
  \item \textbf{True}
  \item \textbf{False}
  \item \textbf{True}
  \item \textbf{True}
  \item \textbf{True}
  \item \textbf{False}
  \item \textbf{True}
  \item \textbf{True}
  \item \textbf{True}
  \item \textbf{False}
  \item \textbf{False}

\end{enumerate}


\end{document}
